\documentclass[11pt,reqno]{amsart}

\usepackage[T1]{fontenc}
\usepackage[utf8]{inputenc}
\usepackage{lmodern}
\usepackage{microtype}
\usepackage{geometry}
\usepackage{amsmath,amssymb,mathtools}
\usepackage{mathrsfs}
\usepackage{booktabs}
\usepackage{longtable}
\usepackage{array}
\usepackage{enumitem}
\usepackage{xcolor}
\usepackage{hyperref}
\usepackage[nameinlink,capitalize,noabbrev]{cleveref}

\geometry{margin=1.08in}
\hypersetup{
  colorlinks=true,
  linkcolor=blue!55!black,
  citecolor=green!35!black,
  urlcolor=blue!65!black,
  pdftitle={Presentations of absolute Galois groups of dyadic local fields},
  pdfauthor={David Roe}
}
\setlist{itemsep=0.25em,topsep=0.45em}
\setcounter{tocdepth}{2}
\allowdisplaybreaks

\newtheorem{theorem}{Theorem}[section]
\newtheorem{proposition}[theorem]{Proposition}
\newtheorem{lemma}[theorem]{Lemma}
\newtheorem{corollary}[theorem]{Corollary}
\newtheorem{conjecture}[theorem]{Conjecture}
\theoremstyle{definition}
\newtheorem{definition}[theorem]{Definition}
\newtheorem{construction}[theorem]{Construction}
\newtheorem{example}[theorem]{Example}
\theoremstyle{remark}
\newtheorem{remark}[theorem]{Remark}
\newtheorem{warning}[theorem]{Warning}

\numberwithin{equation}{section}

\newcommand{\Q}{\mathbf Q}
\newcommand{\Z}{\mathbf Z}
\newcommand{\F}{\mathbf F}
\newcommand{\G}{\mathbf G}
\newcommand{\C}{\mathbf C}
\newcommand{\A}{\mathbf A}
\newcommand{\ur}{\mathrm{ur}}
\newcommand{\tame}{\mathrm{t}}
\newcommand{\pc}{\mathrm{pc}}
\newcommand{\sgn}{\mathrm{sgn}}
\newcommand{\ram}{\mathrm{ram}}
\newcommand{\unr}{\mathrm{unr}}
\newcommand{\prof}{\mathrm{prof}}
\newcommand{\Gal}{\operatorname{Gal}}
\newcommand{\Hom}{\operatorname{Hom}}
\newcommand{\Sur}{\operatorname{Sur}}
\newcommand{\Aut}{\operatorname{Aut}}
\newcommand{\im}{\operatorname{im}}
\newcommand{\coker}{\operatorname{coker}}
\newcommand{\Nm}{\operatorname{Nm}}
\newcommand{\inv}{\operatorname{inv}}
\newcommand{\Tr}{\operatorname{Tr}}
\newcommand{\rad}{\operatorname{rad}}
\newcommand{\Frat}{\operatorname{Frat}}
\newcommand{\eps}{\varepsilon}
\newcommand{\whZ}{\widehat{\Z}}
\newcommand{\simeqto}{\xrightarrow{\sim}}
\newcommand{\dd}{\mathop{}\!\mathrm d}
\newcommand{\Gsum}{\mathcal G}
\newcommand{\cH}{\mathcal H}
\newcommand{\cD}{\mathcal D}
\newcommand{\cY}{\mathcal Y}
\newcommand{\cJ}{\mathcal J}

\DeclareMathOperator{\rank}{rank}
\DeclareMathOperator{\Arf}{Arf}

\title[Dyadic absolute Galois presentations]
{Presentations of Absolute Galois Groups\\of Dyadic Local Fields}
\author{David Roe}
\address{Department of Mathematics, Massachusetts Institute of Technology,
Cambridge, MA 02139, USA}
\email{roed@mit.edu}
\date{\today}

\begin{document}

\begin{abstract}
Let \(K/\Q_2\) be finite.  We give a uniform family of admissible
profinite presentations of the absolute Galois group \(G_K\).  Diekert's
presentation covers the case in which \(K(i)/K\) is unramified.  In the
ramified case, the presentation has \(n+3\) generators, where
\(n=[K:\Q_2]\), a tame relation
\(\tau^\sigma=\tau^{q_K}\), and one explicit wild relation selected from
three cyclotomic types \(L\), \(M_\alpha\), and \(N_\alpha\).

The proof separates into four interfaces.  First, the tame quotient and
the fully marked maximal pro-\(2\) quotient are identified over the full
unramified character.  Second, finite elementary lifting problems are
controlled by the evaluated two-relator Fox complex.  Third, a
Fox--Heisenberg Stokes construction gives the candidate the same
two-dimensional duality package as local Galois cohomology.  Fourth, the
candidate determinant obstruction has the same Gauss sum as the local
one.  The local sign for a nontrivial self-dual simple tame module \(V\)
is \((-1)^n\) when \(V\) is unramified and \(+1\) when \(V\) is ramified.
Finite Fourier inversion and exact-image induction then identify all
finite epimorphism counts, and profinite reconstruction gives
\(G_K\).

This document consolidates a larger collection of component proof notes.
It is a working mathematical draft, not yet a peer-reviewed or
machine-checked paper; the remaining bibliographic and formalization
tasks are recorded explicitly.
\end{abstract}

\maketitle
\tableofcontents

\section{Introduction}

For odd primes \(p\), Jannsen and Wingberg gave presentations of the
absolute Galois groups of \(p\)-adic fields
\cite{JannsenWingberg1982,NSW}.  Residue characteristic two is exceptional:
the usual odd-prime formula contains exponents that do not specialize
integrally at \(p=2\).  Zel'venski\u{\i} described maximal pro-\(2\)
extensions in odd degree \cite{Zelvenskii1972}, and Diekert treated a
substantial class of dyadic absolute Galois groups
\cite{Diekert1984}.  More recently, Roe and Turturean gave an explicit
four-generator presentation of \(G_{\Q_2}\)
\cite{RoeTurturean2026}.  The purpose of the present paper is to extend
their method to every finite extension of \(\Q_2\).

Write
\[
 n=[K:\Q_2],\qquad q_K=\#k_K=2^f,\qquad D_K=G_K(2),
\]
and normalize the geometric unramified character by
\[
 \nu_{\ur}:D_K\twoheadrightarrow\Z_2,\qquad
 \nu_{\ur}(\Phi_{\mathrm{geom}})=1.
\]
Let
\[
 \omega_2=(1\text{ in }\Z_2,\;0\text{ in }\Z_\ell
                    \text{ for every odd }\ell)\in\whZ.
\]
Thus \(\omega_2^2=\omega_2\).  Throughout,
\[
 x^g=g^{-1}xg,\qquad [x,y]=x^{-1}y^{-1}xy.
\]

\subsection{Admissible presentations}

The phrase ``presentation'' below always includes a pro-\(2\) condition.
Given a profinite word \(R\) in
\(\sigma,\tau,x_0,\ldots,x_n\), define
\begin{equation}\label{eq:candidate}
 \Gamma_R=
 \left\langle
 \sigma,\tau,x_0,\ldots,x_n
 \ \middle|\
 \tau^\sigma=\tau^{q_K},\quad R=1,\quad
 \overline{\langle\!\langle x_0,\ldots,x_n\rangle\!\rangle}
 \text{ is pro-}2
 \right\rangle_{\prof}.
\end{equation}
Concretely, \(\Gamma_R\) is the inverse limit of finite quotients in
which the displayed relations hold and the images of all \(x_i\) lie in
the \(2\)-core.  A bare two-relator profinite quotient is not an
equivalent definition.

\begin{theorem}[Main theorem]\label{thm:main}
Let \(K/\Q_2\) be finite.
\begin{enumerate}
\item If \(K(i)/K\) is unramified, \(G_K\) is given by Diekert's
presentation \cite{Diekert1984}.
\item If \(K(i)/K\) is ramified, there is an explicit word \(R_K\),
given in \cref{sec:relators}, such that
\[
 \Gamma_{R_K}\cong G_K.
\]
The word belongs to exactly one of the families \(L\), \(M_\alpha\), or
\(N_\alpha\).
\end{enumerate}
Consequently every dyadic local field has an explicit admissible
profinite presentation.
\end{theorem}

The theorem is not obtained from the abstract Demu\v{s}kin type alone.
The full unramified marking interacts with the cyclotomic orientation,
and that interaction splits the product-type \(M_\alpha\) case into
three word families.

\subsection{Structure of the proof}

The proof has four word-independent interfaces and one word-specific
calculation.
\begin{enumerate}
\item A relative Labute theorem produces the fully marked maximal
pro-\(2\) boundary, preserving both the cyclotomic orientation
\(\chi_K\) and the full character \(\nu_{\ur}\).
\item Evaluating the tame and wild relators in elementary extensions
gives a functorial three-term Fox complex
\[
 0\longrightarrow A\longrightarrow A^{n+3}
   \longrightarrow A^2\longrightarrow0.
\]
\item Finite-word Stokes identities and scalar normalization make this
complex naturally self-dual in dimension two, including nonsplit
coefficient modules.
\item The normalized determinant obstruction is a nonsingular
quadratic form.  Its local Gauss sign is computed uniformly for all
dyadic \(K\).
\item For each displayed word, exact Fox and Hessian calculations verify
the preceding interfaces.  Boundary-framed exact-image induction,
finite Fourier inversion, and profinite reconstruction then finish the
proof.
\end{enumerate}

\section{The marked arithmetic datum}\label{sec:marked-data}

Let
\[
 \chi=\chi_K:D_K\longrightarrow\Z_2^\times
\]
be the cyclotomic orientation.  Put
\begin{equation}\label{eq:CI}
 C=\chi(D_K),\qquad I=\chi(\ker\nu_{\ur}),\qquad
 A=\nu_{\ur}(\ker\chi)=2^r\Z_2.
\end{equation}
The characters induce a canonical quotient map
\begin{equation}\label{eq:lambda}
 \lambda:C\twoheadrightarrow\Z/2^r\Z,\qquad
 \lambda(\chi(g))=\nu_{\ur}(g)\pmod{2^r}.
\end{equation}
If \(\sigma\) is a geometric Frobenius lift, retain also
\begin{equation}\label{eq:gamma}
 \gamma=\chi(\sigma)I\in C/I,\qquad \lambda(\gamma)=1.
\end{equation}

\begin{warning}
The abstract image \(C\), or equivalently the abstract Labute type, does
not determine the marked pair \((D_K,\nu_{\ur})\).  The necessary datum
is \((C,I,\lambda,\gamma)\).  In particular, the three quadratic fields
\(\Q_2(\sqrt5)\), \(\Q_2(\sqrt{10})\), and
\(\Q_2(\sqrt{-10})\) have the same abstract type \(M_2\) but different
marked Frobenius systems.
\end{warning}

The residue cardinality \(q_K\) is different from the Demu\v{s}kin
invariant.  In the branches considered below the latter is \(2\);
\(q_K\) occurs in the tame relation only.

\subsection{Marked Demu\v{s}kin cores}

For even types, use Labute variables \(a,b,c,d\) and
\[
\begin{aligned}
 r_M&=a^2[a,b]c^{2^\alpha}[c,d]\prod_j[u_j,v_j],\\
 r_N&=a^{2+2^\alpha}[a,b][c,d]\prod_j[u_j,v_j].
\end{aligned}
\]
Put \(m=2^{\alpha-1}\).

\paragraph{Type \(L\).}
When \(n\) is odd, \(r=0\), \(I=C\), and the marked core is
\begin{equation}\label{eq:Lcore}
 P_{L,n}=x_0^{\sigma^2}x_0[x_1,\sigma]
 \prod_{j=1}^{(n-1)/2}[x_{2j},x_{2j+1}].
\end{equation}

\paragraph{Type \(M_\alpha\), procyclic Frobenius.}
Write
\[
 u=(1-2^\alpha)^{-1},\qquad
 \lambda(-1)=\epsilon2^{r-1},\qquad
 \lambda(u)=\eta\pmod{2^r}.
\]
If \(\eta\) is odd, put \(\rho=\eta^{-1}\) and
\[
\begin{aligned}
 \sigma&=d^\rho,& x_0&=c^{-m}a^{-1},\\
 x_1&=b\sigma^{-\epsilon2^{r-1}},&
 x_2&=c\sigma^{-2^r}.
\end{aligned}
\]
Equivalently, with
\[
\begin{aligned}
 A&=x_0^{-1}C_0^{-m},&
 B&=x_1\sigma^{\epsilon2^{r-1}},\\
 C_0&=x_2\sigma^{2^r},&
 D&=\sigma^\eta,
\end{aligned}
\]
the marked rank-four core is
\begin{equation}\label{eq:Mpc-core}
 P_{M,\pc}=A^2[A,B]C_0^{2^\alpha}[C_0,D].
\end{equation}

\paragraph{Type \(M_\alpha\), sign Frobenius.}
If the procyclic coordinate is even, surjectivity of \(\lambda\) forces
\[
 r=1,\qquad \epsilon=1.
\]
Use
\[
 \sigma=b,\qquad x_0=c^{-m}a^{-1},\qquad
 x_1=c\sigma^{-2},\qquad x_2=\sigma^{-\eta}d.
\]
With
\[
 A=x_0^{-1}C_0^{-m},\quad B=\sigma,\quad
 C_0=x_1\sigma^2,\quad D=\sigma^\eta x_2,
\]
the core again has the form
\begin{equation}\label{eq:Msign-core}
 P_{M,\sgn}=A^2[A,B]C_0^{2^\alpha}[C_0,D].
\end{equation}

\paragraph{Type \(N_\alpha\).}
For \(r=0\), the core is
\begin{equation}\label{eq:Ncompact-core}
 P_{N,0}=x_0^{2+2^\alpha}[x_0,x_1][\sigma,x_2].
\end{equation}
For \(r\ge1\), let
\[
 v=-(1+2^\alpha)^{-1},\qquad
 \eta=\lambda(v)\in\Z_2^\times,\qquad \rho=\eta^{-1}.
\]
Take
\[
 \sigma=b^\rho,\qquad x_0=a,\qquad
 x_1=c\sigma^{-2^r},\qquad x_2=d.
\]
Then
\begin{equation}\label{eq:Npc-core}
 P_{N,r,\eta}
 =x_0^{2+2^\alpha}[x_0,\sigma^\eta]
  [x_1\sigma^{2^r},x_2].
\end{equation}
Append \((n-2)/2\) hyperbolic handles to every even core.

\subsection{Uniform relative Labute theorem}

\begin{theorem}[Two-character Labute normalization]\label{thm:labute}
Let \(F\) be a finitely generated free pro-\(2\) group carrying
characters \(\chi\) and \(\nu\), and let \(P\) be one of the marked cores
above, stabilized by hyperbolic handles.  Suppose that \(r_0\in\Frat(F)\)
is a Demu\v{s}kin relator with
\begin{enumerate}
\item the same Labute type as \(P\);
\item canonical orientation \(\chi\);
\item the same abelian relation vector and quadratic initial form as
\(P\); and
\item \(\nu(r_0)=0\).
\end{enumerate}
Then an automorphism \(\Theta\in\Aut(F)\) satisfies
\[
 \Theta\!\left(\overline{\langle\!\langle r_0\rangle\!\rangle}\right)
 =\overline{\langle\!\langle P\rangle\!\rangle},\qquad
 \chi\Theta=\chi,\qquad \nu\Theta=\nu.
\]
\end{theorem}

\begin{proof}[Proof outline]
Local reciprocity and Smith reduction put the two characters in one of
the rows above.  The identity \(z\cup z=z\cup\kappa\), followed by Witt
cancellation, puts the mod-\(2\) cup form into the displayed rank-three
or rank-four core plus hyperbolic handles.

The remaining issue is integral and all-degree.  In the restricted Lie
algebra attached to the Zassenhaus filtration, the degree-\(j\)
difference between \(r_0\) and \(P\) lies in the image of the Nielsen
differential after adjoining the three canonical orientation
functionals.  Compatibility with both characters annihilates precisely
the exceptional three-dimensional cokernel.  Choose each correction in
\(\ker\chi\cap\ker\nu\), eliminate handle terms first, and iterate in
degrees \(j=3,4,\ldots\).  The correction degrees tend to infinity, so
the product converges to a continuous free pro-\(2\) automorphism.  Every
factor preserves both characters.
\end{proof}

\begin{corollary}[Marked maximal pro-\(2\) quotient]\label{cor:marked-D}
If \(K(i)/K\) is ramified, the appropriate stabilized core \(P_K\)
presents \(D_K\) as an oriented, unramified-marked pro-\(2\) group:
\[
 (D_K,\chi_K,\nu_{\ur})\cong
 \left(F/\overline{\langle\!\langle P_K\rangle\!\rangle},
       \chi_{K}^{\mathrm{std}},
       \nu_{K}^{\mathrm{std}}\right),
\]
where
\[
 \nu_{K}^{\mathrm{std}}(\sigma)=1,\qquad
 \nu_{K}^{\mathrm{std}}(x_i)=0.
\]
\end{corollary}

\section{The field-independent source theorem}\label{sec:source}

Let
\[
 T_q=\langle\sigma,\tau\mid\tau^\sigma=\tau^q\rangle_{\prof},
\qquad
 \nu_{\tame}(\sigma)=1,\quad \nu_{\tame}(\tau)=0,
\]
and define the common boundary
\begin{equation}\label{eq:boundary}
 \partial_K=T_{q_K}\times_{\Z_2}D_K.
\end{equation}
For \(S=\Gamma_R\) or \(G_K\), write \(W_S=O_2(S)\).

\subsection{The four interfaces}

\begin{definition}[B1: fully marked boundary]\label{def:B1}
The source \(S\) satisfies \emph{B1} if it has compatible epimorphisms
\[
 S\twoheadrightarrow T_{q_K},\qquad S\twoheadrightarrow D_K
\]
such that \(S/W_S\cong T_{q_K}\), the second map is the maximal
pro-\(2\) quotient, the induced characters to \(\Z_2\) agree, the
resulting map \(S\twoheadrightarrow\partial_K\) is surjective, and the
candidate marking is
\[
 \nu_{\ur}(\sigma)=1,\qquad \nu_{\ur}(x_i)=0.
\]
\end{definition}

\begin{definition}[B2: exact finite lifting semantics]\label{def:B2}
Fix a finite lower exact-image map \(\rho:S\to C\) and an elementary
\(\F_2[C]\)-module \(A\).  The candidate satisfies \emph{B2} if
evaluation of the two relators gives a functorial complex
\begin{equation}\label{eq:fox-complex}
 C^\bullet_{\Gamma_R,\rho}(A):
 0\longrightarrow A\xrightarrow{d^0}A^{n+3}
  \xrightarrow{d_R^1}A^2\longrightarrow0
\end{equation}
with the following literal lifting interpretation:
\begin{enumerate}
\item changing generator lifts by \(a\in A^{n+3}\) changes the relation
defect by \(d_R^1(a)\);
\item the defect class in \(\coker d_R^1\) vanishes exactly when a lift
exists;
\item when nonempty, the raw lift set is a torsor under
\(\ker d_R^1\); and
\item the construction is natural, boundary-compatible, and exact for
short exact coefficient sequences.
\end{enumerate}
\end{definition}

\begin{definition}[B3: linear duality]\label{def:B3}
The candidate satisfies \emph{B3} if there is a natural
quasi-isomorphism
\begin{equation}\label{eq:duality-qiso}
 C^\bullet_{\Gamma_R,\rho}(A)\simeq
 \Hom_{\F_2}\!\left(
 C^\bullet_{\Gamma_R,\rho}(A^\vee),\F_2\right)[-2],
\end{equation}
whose cohomological pairings agree with local Tate duality.  In
particular,
\[
 H^2_{\Gamma_R,\rho}(A)^\vee\cong(A^\vee)^C,
\]
connecting homomorphisms are adjoint, and for \(A=\F_2\) the
cup--Bockstein form agrees with the local Hilbert form under the marked
maximal pro-\(2\) identification.
\end{definition}

\begin{definition}[B4: quadratic determinant]\label{def:B4}
Let \(V\) be a nontrivial self-dual simple tame
\(\F_2[C]\)-module with invariant nonsingular quadratic form \(q_V\).
Let \(Q^0_{S,\rho}\) be the zero-section-normalized determinant
obstruction on \(H^1_{S,\rho}(V)\).  The candidate satisfies \emph{B4}
if
\begin{equation}\label{eq:B4}
 \sum_{c\in H^1_{\Gamma_R,\rho}(V)}
  (-1)^{Q^0_{\Gamma_R,\rho}(c)}
 =
 \sum_{c\in H^1_{G_K,\rho}(V)}
  (-1)^{Q^0_{G_K,\rho}(c)}.
\end{equation}
\end{definition}

\begin{theorem}[Source-interface theorem]\label{thm:source}
Suppose that \(\Gamma_R\) and \(G_K\) have the common boundary
\eqref{eq:boundary} and satisfy B1--B4.  Then every boundary-framed
finite target and every exact-image stratum has the same number of
epimorphisms from \(\Gamma_R\) and \(G_K\).  Consequently,
\[
 \#\Sur(\Gamma_R,G)=\#\Sur(G_K,G)
 \quad\text{for every finite group }G,
\]
and
\[
 \Gamma_R\cong G_K.
\]
\end{theorem}

\begin{proof}[Proof architecture]
Fix a boundary-framed target \(Y\) with normal \(2\)-subgroup \(L\).
The argument is by strong induction on \(|L|\).

If every simple factor of \(L/\Frat(L)\) is trivial, B1 and the scalar
part of B3 identify the terminal lift counts.  Otherwise choose a
minimal non-scalar simple head \(V\).  B2 identifies each nonempty
linear lift fibre with a \(Z^1(V)\)-torsor.  B3 gives the same
dimensions, obstruction pairings, and adjoint boundary maps on the two
sources.  For non-self-dual \(V\), this already equates the required
fibres.

For self-dual \(V\), scalar central pushouts introduce a quadratic
determinant obstruction.  B4 identifies its Fourier transform, hence
the zero and nonzero fibres.  Central Frattini covers and
inclusion--exclusion over exact image subgroups reduce all remaining
terms to smaller normal \(2\)-subgroups or proper exact-image strata.
The induction therefore closes.

Summing boundary frames gives equality of unframed epimorphism counts.
The candidate is topologically finitely generated; the one-sided
profinite reconstruction theorem then identifies it with \(G_K\)
\cite{DixonEtAl1982,RibesZalesskii,FriedJarden}.
\end{proof}

\subsection{Numerical consequences}

Let \(d=\dim_{\F_2}A\).  The Euler characteristic of
\eqref{eq:fox-complex} is \(-nd\).  B3 gives
\[
 h^0(A)=\dim A^C,\qquad h^2(A)=\dim(A^\vee)^C,
\]
and hence
\begin{align}
 \dim H^1(A)
 &=nd+\dim A^C+\dim(A^\vee)^C,\label{eq:h1dim}\\
 \dim Z^1(A)
 &=(n+1)d+\dim(A^\vee)^C.\label{eq:z1dim}
\end{align}
Thus
\[
 \dim H^1(\F_2)=n+2,\qquad \dim H^2(\F_2)=1,
\]
whereas for a nontrivial simple \(V\),
\[
 \dim H^1(V)=n\dim V,\qquad
 \dim Z^1(V)=(n+1)\dim V.
\]

\section{The local determinant sign}\label{sec:local-gauss}

Let \(V\) be a nontrivial finite simple continuous
\(\F_2[G_K]\)-module.  A normal pro-\(2\) subgroup acts trivially on
\(V\), so its action factors through a finite tame image \(H\).  Suppose
\(V\) is self-dual and admits an \(H\)-invariant nonsingular quadratic
form \(q:V\to\F_2\).  Let
\[
 \kappa_q^0\in H^2(V\rtimes H,\F_2)
\]
be the equivariant extraspecial class normalized to restrict trivially
to the zero section.  For a cocycle \(b\), define
\begin{equation}\label{eq:local-Q}
 Q^0_{K,V}([b])
 =\inv_K\!\left((b,\rho)^*\kappa_q^0\right).
\end{equation}

\begin{theorem}[Arbitrary-dyadic Gauss sign]\label{thm:local-gauss}
The form \(Q^0_{K,V}\) is well-defined and nonsingular on
\(H^1(K,V)\), whose dimension is \(n\dim V\).  Moreover,
\begin{equation}\label{eq:local-gauss}
 \Gsum(Q^0_{K,V})
 :=\sum_{x\in H^1(K,V)}(-1)^{Q^0_{K,V}(x)}
 =\eps_K(V)\,2^{n\dim(V)/2},
\end{equation}
where
\begin{equation}\label{eq:local-sign}
 \eps_K(V)=
 \begin{cases}
 (-1)^n,&V\text{ is unramified},\\
 +1,&V\text{ is ramified}.
 \end{cases}
\end{equation}
Equivalently,
\[
 \#(Q^0_{K,V})^{-1}(0)
 =2^{n\dim V-1}+\eps_K(V)2^{n\dim V/2-1}.
\]
\end{theorem}

\begin{proof}
Local Euler--Poincar\'e and Tate duality give the dimension and
nonsingularity \cite[Chapter VII]{NSW}.  It remains to compute the Arf
invariant.

Suppose first that \(V\) is ramified.  Pass to the faithful tame
splitting field and use the dyadic square-class filtration.  The middle
layers are parity-free: when the inertia order is even the relevant
isotypic layer is zero, while for odd inertia order a nonzero middle
layer would force trivial inertia.  Exactness of the faithful isotypic
functor follows from projectivity over a Sylow subgroup.  The deep unit
part therefore gives a totally singular subspace of dimension
\(n\dim(V)/2\).  It is a Lagrangian, so the form is split and its Gauss
sign is \(+1\).  The normalized orbit calculation uses the
Evens--Kahn formula \cite{Evens1963,Kahn1984,Kozlowski1984}.

If \(V\) is unramified, the coefficient field of its simple endomorphism
algebra carries the adjoint involution induced by self-duality.
The local form becomes a nonsingular Hermitian form of rank \(n\).  A
single Hermitian line has minus Arf sign, and orthogonal sum multiplies
Gauss signs.  Thus rank \(n\) contributes \((-1)^n\).
\end{proof}

\begin{remark}
For even \(n\), both ramified and unramified simples have plus sign.
For odd \(n\), the unramified sign is negative and the ramified sign is
positive.  This theorem supplies the target sign; every candidate word
must still reproduce it by an independent Hessian calculation.
\end{remark}

\section{Explicit \texorpdfstring{ramified-\(i\)}{ramified-i} relators}
\label{sec:relators}

All groups in this section are admissible in the sense of
\eqref{eq:candidate}.  Put
\[
 \sigma_2=\sigma^{\omega_2},\qquad
 \delta_i=(x_i\tau)^{\omega_2}x_i^{-1}.
\]
For the even-degree families write
\[
 n=2+2h,\qquad
 H_h=\prod_{j=1}^{h}[x_{2j+1},x_{2j+2}].
\]
Every product below is taken in the displayed order.

\subsection{Odd degree: type \texorpdfstring{\(L\)}{L}}

Let \(n=2m+1\).  Put
\[
 u_i=(x_i\tau)^{\omega_2}\quad(i=0,1),
\]
\[
\begin{aligned}
 d_0&=u_0x_0^{-1},&
 z_0&=x_0^{\sigma_2},&
 c_0&=[d_0,z_0],\\
 g_0&=\sigma_2^2,&
 d_g&=d_0^{g_0},&
 h_c&=[d_g,d_0],\\
 h_0&=x_0^{g_0}x_0\,d_gd_0d_0^2h_c.
\end{aligned}
\]
Define
\begin{equation}\label{eq:Lword}
 R_{L,K}
 =h_0u_1^{-1}x_1^\sigma c_0
  \prod_{j=1}^{m}[x_{2j},x_{2j+1}].
\end{equation}

\begin{theorem}\label{thm:Lword}
If \(n\) is odd, then
\[
 \Gamma_{R_{L,K}}\cong G_K.
\]
\end{theorem}

\begin{remark}
The cyclotomic core uses \(g_0=\sigma_2^2\), not
\(\sigma_2^{q_K}\).  The residue cardinality changes the tame relation
and its Fox norm polynomial, not the cyclotomic depth.
\end{remark}

\subsection{Type \texorpdfstring{\(N_\alpha\)}{N-alpha}}

Put \(p_\alpha=2+2^\alpha\).

\paragraph{Compact row.}
If \(r=0\), define
\begin{equation}\label{eq:Ncompact-word}
 R_{N,\alpha,0}
 =x_0^{p_\alpha}[x_0,x_1]\,
  x_2^{-\sigma}(x_2\tau)^{\omega_2}H_h.
\end{equation}

\paragraph{Nontrivial Frobenius quotient.}
Suppose \(r\ge1\).  Let
\[
 \eta=\lambda\!\left(-(1+2^\alpha)^{-1}\right)\in\Z_2^\times
\]
and let \(\widehat\eta\in\whZ^\times\) have \(2\)-component \(\eta\)
and odd-primary components \(1\).  Put
\[
 g=x_1\sigma^{2^r},\qquad
 D_{r,\eta}
 =\delta_0^{\sigma^{\widehat\eta}}\,
  \delta_0^{\sigma^{-2^r}}\,
  \delta_0^{\sigma^{\widehat\eta-2^r}},
\qquad
 E_{r,\eta}=[D_{r,\eta},x_1].
\]
Define
\begin{equation}\label{eq:Npc-word}
 R_{N,\alpha,r,\eta}
 =x_0^{p_\alpha}[x_0,\sigma^{\widehat\eta}]\,
 x_2^{-g}(x_2\tau)^{\omega_2}E_{r,\eta}H_h.
\end{equation}

\begin{theorem}\label{thm:Nwords}
For every ramified-\(i\) field of type \(N_\alpha\),
\[
 \Gamma_{R_{N,\alpha,0}}\cong G_K\quad(r=0),
\qquad
 \Gamma_{R_{N,\alpha,r,\eta}}\cong G_K\quad(r\ge1).
\]
\end{theorem}

The correction \(E_{r,\eta}\) is invisible in the tame quotient, in the
maximal pro-\(2\) quotient, and at first Fox order.  At second order it is
essential.  On ramified simples, writing
\[
 A=S^{\widehat\eta},\qquad B=S^{2^r},
\]
it changes the two cross operators to
\begin{equation}\label{eq:Ncross}
 M_c=A,\qquad L_c=A^{-1}.
\end{equation}
A two-dimensional \(S_3\)-module already detects the radical in the
uncorrected word.

\subsection{Type \texorpdfstring{\(M_\alpha\)}{M-alpha}}

Put
\[
 m=2^{\alpha-1},\qquad u=(1-2^\alpha)^{-1}.
\]
Exactly three families occur.

\subsubsection{Compact row}

If \(r=0\), put
\[
 A_0=x_0^{-1}\sigma_2^{-m},\qquad
 J_2=x_2^{-\sigma}(x_2\tau)^{\omega_2},
\]
\[
 E_m^{\mathrm{rev}}
 =\delta_1^{\sigma_2^{2m}}
  \delta_1^{\sigma_2^m}
  \delta_0^{\sigma_2^m}
  \delta_0.
\]
The word is
\begin{equation}\label{eq:Mcompact-word}
 R_{M,0}
 =A_0^2[A_0,x_1]\sigma_2^{2m}
 J_2E_m^{\mathrm{rev}}H_h.
\end{equation}
The reversal in \(E_m^{\mathrm{rev}}\) is essential; the forward order is
singular on a fifth-root orbit.

\subsubsection{Procyclic-Frobenius row}

Assume \(r\ge1\) and
\[
 \eta=\lambda(u)\in\Z_2^\times.
\]
Put
\[
 s=2^r,\qquad p=\epsilon2^{r-1},
\]
and let \(\widehat\eta\in\whZ^\times\) have \(2\)-component \(\eta\)
and odd-primary components \(1\).  Define
\[
 C_0=x_2\sigma_2^s,\quad
 A=x_0^{-1}C_0^{-m},\quad
 B=x_1\sigma_2^p,\quad
 D=\sigma^{\widehat\eta},
\]
\begin{align}
 E_{01}^{\pc}
 &=
 \delta_1^{\sigma_2^{p+2sm}}
 \delta_1^{\sigma_2^{p+sm}}
 \delta_0^{\sigma_2^{sm+p}}
 \delta_0,\label{eq:E01pc}\\
 E_2^{\pc}
 &=
 \delta_2^{\sigma_2^s}
 \prod_{j=1}^{m}
 \left(
 \delta_2^{\sigma_2^{s(m+j)}}
 \delta_2^{\sigma_2^{p+s(m+j)}}
 \right).\label{eq:E2pc}
\end{align}
The product in \eqref{eq:E2pc} is in increasing \(j\).  Put
\[
 R_{\mathrm{lin}}^{\pc}
 =A^2[A,B]C_0^{2^\alpha}[C_0,D]\,
 E_{01}^{\pc}E_2^{\pc}.
\]
With \(D_i=\delta_i\), define
\[
 \widehat C=\sigma_2^s,\quad
 \widehat A=D_0^{-1}\widehat C^{-m},\quad
 \widehat B=D_1\sigma_2^p,\quad
 \widehat D=\sigma^{\widehat\eta},
\]
\[
 \widehat R^{\pc}
 =
 \widehat A^2[\widehat A,\widehat B]
 \widehat C^{2^\alpha}[\widehat C,\widehat D]\,
 E_{01}^{\pc}.
\]
The final word is
\begin{equation}\label{eq:Mpc-word}
 R_{M,\pc}
 =R_{\mathrm{lin}}^{\pc}\widehat R^{\pc}
 D_0^2[D_0,D_1]H_h.
\end{equation}

\subsubsection{Sign-Frobenius row}

If \(\eta=\lambda(u)\) is even, then
\[
 r=1,\qquad\epsilon=1.
\]
Use the left-adapted basis and put
\[
 C_0=x_1\sigma_2^2,\quad
 A=x_0^{-1}C_0^{-m},\quad
 D=\sigma_2^\eta x_2,\quad
 F_2=[\sigma,x_2]^{-1}x_2^{-\sigma}
       (x_2\tau)^{\omega_2}.
\]
Define
\[
 E_0^{\sgn}=\delta_0^{\,\sigma\sigma_2^{2m}}\delta_0,
\]
\begin{equation}\label{eq:E1sign}
 E_1^{\sgn}
 =
 \delta_1^{\sigma_2^{2+\eta}}
 \delta_1^{\sigma_2^2}
 \prod_{j=m}^{1}
 \left(
 \delta_1^{\,\sigma\sigma_2^{2(m+j)}}
 \delta_1^{\,\sigma_2^{2(m+j)}}
 \right),
\end{equation}
where the product is in decreasing \(j\), and put
\[
 E_{01}^{\sgn}=E_1^{\sgn}E_0^{\sgn},
\]
\[
 R_{\mathrm{lin}}^{\sgn}
 =A^2[A,\sigma]C_0^{2^\alpha}[C_0,D]\,
 F_2E_{01}^{\sgn}.
\]
With \(D_i=\delta_i\), define
\[
 \widehat C=D_1\sigma_2^2,\quad
 \widehat A=D_0^{-1}\widehat C^{-m},\quad
 \widehat D=\sigma_2^\eta,
\]
\[
 \widehat R^{\sgn}
 =\widehat A^2[\widehat A,\sigma]
 \widehat C^{2^\alpha}[\widehat C,\widehat D]\,
 E_{01}^{\sgn}.
\]
The final word is
\begin{equation}\label{eq:Msign-word}
 R_{M,\sgn}
 =R_{\mathrm{lin}}^{\sgn}\widehat R^{\sgn}
 D_0^2[D_0,D_1]H_h.
\end{equation}

\begin{theorem}\label{thm:Mwords}
For every ramified-\(i\) field of type \(M_\alpha\), exactly one of
\eqref{eq:Mcompact-word}, \eqref{eq:Mpc-word}, and
\eqref{eq:Msign-word} applies, and the corresponding admissible group is
isomorphic to \(G_K\).
\end{theorem}

\begin{remark}[Self-replication]
On a ramified simple module, the auxiliary word \(\widehat R\) has zero
first Fox derivative and exactly reproduces the entire raw extraspecial
determinant, including the \(T\)-dependent central terms.  The two copies
cancel in characteristic two, and \(D_0^2[D_0,D_1]\) leaves the canonical
plus form
\[
 Q_+(c_0,c_1)=q(c_0)+b_q(c_0,c_1).
\]
This is a word identity, not an interpolation from finite inertia
orders.
\end{remark}

\section{Verification of the interfaces}\label{sec:verification}

We now explain why the words in \cref{sec:relators} satisfy B1--B4.
The calculations are uniform once the rank-three or rank-four core has
been fixed.

\subsection{Boundary and lifting semantics}

Kill every \(x_i\).  Each \(\delta_i\), auxiliary correction, and handle
becomes trivial; the wild relation contributes no new tame relation.
Thus the quotient is \(T_{q_K}\).  Kill \(\tau\) and replace
\(\sigma_2\) by \(\sigma\).  Each correction becomes trivial and the
wild word specializes to the marked core of
\cref{sec:marked-data}.  By \cref{cor:marked-D}, this is \(D_K\) with
the full unramified marking.  Goursat's lemma, applied over \(\Z_2\),
gives the surjection to \(\partial_K\).  This proves B1.

For B2, fix an extension
\[
 1\longrightarrow A\longrightarrow B\longrightarrow C\longrightarrow1
\]
and chosen lifts of the \(n+3\) generators.  The tame and wild words
evaluate to a defect in \(A^2\).  Replacing the lifts by
\(a\in A^{n+3}\) changes that defect by the literal evaluated Fox row
\(d_R^1(a)\).  Hence a lift exists exactly when the defect dies in
\(\coker d_R^1\), and all lifts form a torsor under \(\ker d_R^1\).
Every correction used above is either accounted for in the displayed
Fox row or has exact zero first derivative.  The pro-\(2\) condition
adds no linear equation because an elementary kernel already lies in
the marked normal \(2\)-subgroup.

\subsection{Finite-word Stokes and linear duality}

The first- and second-order calculations are made in finite semidirect,
Heisenberg, and extraspecial groups.  If \(u,v\) are words with lower
values \(\bar u,\bar v\), then
\begin{align}
 D(uv)&=D(u)+\bar uD(v),&
 D(u^{-1})&=-\bar u^{-1}D(u),\label{eq:fox-rules}\\
 \beta(uv)&=\beta(u)+\beta(v)
  +D^\vee(u)(\bar uD(v)).\label{eq:heis-rule}
\end{align}
Conjugates, commutators, and profinite powers follow recursively.
Because every coefficient calculation factors through a finite group,
the exponent \(\omega_2\) may be represented by any integer with the
correct residue modulo the group exponent; the derivatives are
independent of that representative.

Summing the two relation defects gives a finite-word Stokes pairing
\[
 C^\bullet(A)\otimes C^\bullet(A^\vee)
 \longrightarrow\F_2[-2].
\]
The individual tame and wild terms are not separately multiplicative:
the cross effects cancel only after both relators are included.  For
each nontrivial simple module, the exact Fox normal form reduces the
chain map to one of the nonsingular core matrices below, plus standard
hyperbolic handles.  For \(A=\F_2\), direct evaluation gives the local
Hilbert cup--Bockstein matrix under the marked isomorphism of B1.

For a short exact coefficient sequence
\[
 0\longrightarrow A'\longrightarrow A\longrightarrow A''\longrightarrow0,
\]
the word complexes form a termwise exact sequence, and the finite-word
Stokes maps give a commuting dual diagram.  The mapping cones are again
termwise exact.  Induction on a composition series shows that the
duality map is a quasi-isomorphism for every elementary coefficient
module, including nonsplit modules.  This proves B3, not merely equality
of dimensions.

\subsection{Candidate quadratic forms}

Let \(V\) be a nontrivial self-dual simple tame module,
\(d=\dim_{\F_2}V\), and let \(b_q\) be the polarization of its invariant
quadratic form.

\paragraph{Odd \(L\).}
The \(n=1\) core is the Roe--Turturean form
\cite{RoeTurturean2026}.  Every additional handle contributes a standard
hyperbolic block.  Thus the unramified sign is multiplied by
\((-1)^{n-1}=+1\), retaining the required minus sign for odd \(n\),
whereas the ramified core remains plus.

\paragraph{\(N_\alpha\).}
For \(r=0\),
\[
 Q(c_0,c_1)=q(c_0)+b_q(c_0,c_1).
\]
For \(r\ge1\), the correction in \eqref{eq:Npc-word} gives
\[
 Q(c_0,c_1)=Q_0(c_0)+b_q(c_1,L_cc_0),
\qquad L_c\in\Aut(V).
\]
The \(c_1\)-subspace is a totally singular Lagrangian.  The rank-four
core therefore has Gauss sum \(2^d\), and the \(h\) handles yield
\[
 \Gsum(Q)=2^{(h+1)d}=2^{nd/2}.
\]

\paragraph{\(M_\alpha\).}
For the compact core, the two simple normal forms give
\[
\begin{array}{c|c}
 \text{inertia projector}&\text{determinant core}\\ \hline
 P=1&q(c_0)+b_q(c_0,c_1)\\
 P=0&q(c_1)+b_q(c_0,c_1).
\end{array}
\]
Both are plus forms.  For the procyclic and sign rows, the
self-replication correction gives the plus form on ramified simples.
The unramified procyclic core is again plus, while the unramified sign
core is hyperbolic.  Appending handles gives
\[
 \Gsum(Q)=2^{nd/2}.
\]

Since all \(M_\alpha\) and \(N_\alpha\) fields have even degree, these
values equal the local sign in \cref{thm:local-gauss}.  This proves B4
for every displayed word.

\begin{corollary}\label{cor:branch-source}
Every word in \cref{thm:Lword,thm:Nwords,thm:Mwords} satisfies B1--B4.
Therefore the corresponding candidate is \(G_K\) by
\cref{thm:source}.
\end{corollary}

\section{Quadratic fields and the formerly missing cases}
\label{sec:quadratic}

The ramified-\(i\) quadratic fields provide useful concrete
specializations.  They also expose errors that are invisible to small
finite-target tests.

\subsection{\texorpdfstring{\(\Q_2(\sqrt{-2})\)}{Q2(sqrt(-2))}}

This is type \(N_2\), with \(I=C=\langle3\rangle\), \(r=0\), and
\(q_K=2\).  The relation is
\[
 x_0^6[x_0,x_1]x_2^{-\sigma}
 (x_2\tau)^{\omega_2}=1.
\]
Thus
\[
\left\langle
\sigma,\tau,x_0,x_1,x_2
\ \middle|\
\tau^\sigma=\tau^2,\quad
x_0^6[x_0,x_1]x_2^{-\sigma}(x_2\tau)^{\omega_2}=1,\quad
\overline{\langle\!\langle x_0,x_1,x_2\rangle\!\rangle}
\text{ pro-}2
\right\rangle_{\prof}
\cong G_{\Q_2(\sqrt{-2})}.
\]

\subsection{\texorpdfstring{\(\Q_2(\sqrt2)\) and
\(\Q_2(\sqrt5)\)}{Q2(sqrt(2)) and Q2(sqrt(5))}}

These are compact \(M_\alpha\) fields.  In
\eqref{eq:Mcompact-word}, take
\[
 m=4\quad(d=2,\ M_3),\qquad
 m=2\quad(d=5,\ M_2).
\]
The reversed factor order in \(E_m^{\mathrm{rev}}\) is necessary.  The
forward order has determinant
\[
 R^{-2}(1+R+R^2+R^3+R^4)
\]
on the ramified mixed operator and is singular on a primitive
fifth-root orbit.

\subsection{\texorpdfstring{\(\Q_2(\sqrt{10})\)}{Q2(sqrt(10))}}

The correct marking in Labute variables is
\[
 \nu(a,b,c,d)=(-4,0,2,1),
\]
not \((0,0,0,1)\).  Put
\[
 \sigma=d,\qquad x_0=c^{-2}a^{-1},\qquad
 x_1=b,\qquad x_2=c\sigma^{-2},
\]
\[
 C_0=x_2\sigma_2^2,\qquad A=x_0^{-1}C_0^{-2},
\]
\[
 E_{10}
 =\delta_1^{\sigma_2^8}
  \delta_1^{\sigma_2^4}
  \delta_0^{\sigma_2^4}
  \delta_0
  \delta_2^{\sigma_2^2}.
\]
The exact specialized word is
\[
 R_{10}=A^2[A,x_1]C_0^4[C_0,\sigma]E_{10}.
\]
It realizes the required inertia orientation image
\(\{\pm1\}(1+8\Z_2)\).

\subsection{\texorpdfstring{\(\Q_2(\sqrt{-10})\)}{Q2(sqrt(-10))}}

The general sign-Frobenius word \eqref{eq:Msign-word} applies.  A useful
field-specific alternative makes the \(T\)-dependent correction
explicit.  Put
\[
 A_-=x_0^{-1}x_1^{-2},
\]
\[
 R_-^{\mathrm{raw}}
 =A_-^2[A_-,\sigma]x_1^4[x_1,\sigma x_2]\,
 [\sigma,x_2]^{-1}x_2^{-\sigma}(x_2\tau)^{\omega_2}.
\]
With \(D_i=\delta_i\), set
\[
 W=D_0D_0^\sigma D_1D_1^\sigma,\qquad
 \Nm_\tau(Z)=(Z\tau)^{\omega_2}\tau^{-\omega_2},\qquad
 \widehat A=D_0^{-1}D_1^{-2},
\]
\[
\widehat R_-=
\widehat A^2[\widehat A,\sigma]D_1^4
[D_1,\sigma W][\sigma,W]^{-1}
W^{-\sigma}\Nm_\tau(W),
\]
\[
 \Theta_-^T=\widehat R_-D_0^2[D_0,D_1],\qquad
 R_-^{\mathrm{new}}=R_-^{\mathrm{raw}}\Theta_-^T.
\]
The factor \(\tau^{-\omega_2}\) removes an unwanted first-order term but
retains the central ramified operator.  The earlier correction
\(D_0^2[D_0,D_1]\) by itself is false: on
\[
 V=\F_{64},\qquad T(x)=\zeta x\ (|\zeta|=9),\qquad S(x)=x^{32},
\]
with
\[
 q(x)=\Tr_{\F_8/\F_2}(x^9),
\]
its mixed Hessian has a two-dimensional radical.  The relative-norm
word above repairs this exact order-nine obstruction.

\section{Assembly of the main theorem}\label{sec:assembly}

\begin{proof}[Proof of \cref{thm:main}]
If \(K(i)/K\) is unramified, invoke Diekert
\cite{Diekert1984}.  Suppose \(K(i)/K\) is ramified.  Labute's
classification \cite{Labute1967}, refined by the marked datum
\((C,I,\lambda,\gamma)\), puts \(D_K\) in exactly one of the following
rows:
\begin{center}
\begin{tabular}{lll}
\toprule
Arithmetic branch & Word & Result\\
\midrule
\(n\) odd, type \(L\) & \eqref{eq:Lword} & \cref{thm:Lword}\\
\(N_\alpha,\ r=0\) & \eqref{eq:Ncompact-word} & \cref{thm:Nwords}\\
\(N_\alpha,\ r\ge1\) & \eqref{eq:Npc-word} & \cref{thm:Nwords}\\
\(M_\alpha,\ r=0\) & \eqref{eq:Mcompact-word} & \cref{thm:Mwords}\\
\(M_\alpha,\ \eta\text{ odd}\) & \eqref{eq:Mpc-word} & \cref{thm:Mwords}\\
\(M_\alpha,\ \eta\text{ even}\) & \eqref{eq:Msign-word} & \cref{thm:Mwords}\\
\bottomrule
\end{tabular}
\end{center}
There is no fourth \(\lambda\)-row: in product type, even \(\eta\)
forces \(r=1,\epsilon=1\).

By \cref{cor:branch-source}, the corresponding candidate satisfies
B1--B4.  The source-interface theorem gives equality of finite
epimorphism counts and hence the desired profinite isomorphism.
\end{proof}

\section{Computational and formal checks}\label{sec:checks}

The proof does not depend on finite testing.  Nevertheless, independent
regressions are useful because the relators are highly sensitive to
factor order and conjugation convention.

\begin{longtable}{p{0.31\textwidth}p{0.61\textwidth}}
\toprule
Check & Result\\
\midrule
\endhead
\(N_2\) class-three word evaluation&
All \(31\) assertions passed.  Independent evaluation paths agreed;
wrong conjugation side, reversed norm order, and even-quotient mutants
were rejected.\\
\addlinespace
Relative Labute restricted Lie algebra&
Degrees \(3\)--\(8\) reproduced the exceptional three-dimensional
cokernel and its annihilation by the three orientation functionals.\\
\addlinespace
General \(N_\alpha\)&
Ramified \(S_3\), unramified \(C_3\), and even-residue-degree checks
gave the predicted corrected ranks and Heisenberg identities.\\
\addlinespace
\(d=-10\) operator&
The old correction fails at inertia order \(9\); the relative-norm
correction is canonical on every tested self-dual odd inertia orbit
through order \(197\).\\
\addlinespace
General \(M_\alpha\)&
The new \(r=2,\epsilon=1\), inertia-order-\(17\) procyclic case has
canonical rank \(16\); ramified sign and unramified sign tests have the
predicted rank-four canonical and hyperbolic forms.\\
\addlinespace
Finite-target oracle&
Direct relator counts agree with an independent induced-wreath oracle
for tested abelian targets, \(S_3\), \(D_8\), and \(A_4\).\\
\addlinespace
Lean prototype&
The interface skeleton compiles with no declared axioms.  It checks
types and imports only; it does not formalize the presentation
theorems.\\
\bottomrule
\end{longtable}

For reference, the validated nonabelian epimorphism counts in the five
quadratic ramified-\(i\) cases are
\[
\begin{array}{c|rrrrr}
 &-2&2&5&10&-10\\ \hline
 S_3&6&6&0&6&6\\
 D_8&1568&1568&1568&1568&1568\\
 A_4&120&120&480&120&120.
\end{array}
\]
The fresh \(A_4\) direct run was compared with previously completed
independent oracle artifacts; a duplicate full oracle recomputation was
stopped after exceeding the allotted runtime without a checkpoint.

\section{Lean formalization plan}\label{sec:lean}

The existing Lean development for \(G_{\Q_2}\) already contains substantial
parts of the finite-group, cohomological, quadratic, and reconstruction
engines.  The general project should therefore proceed by extracting
source-specific leaves into parameterized interfaces, not by rewriting the
finite-target induction from scratch.  This section fixes the intended API
and the order in which it should be implemented.

\subsection{The formal target}

Let \(K\) be represented as a finite intermediate field of
\(\overline{\Q}_2/\Q_2\), and write
\[
 G_K=\operatorname{Gal}(\overline{\Q}_2/K).
\]
For a selected relator family, the final Lean statement should have the
shape
\begin{verbatim}
theorem dyadicPresentation
    (K : IntermediateField Q2 Q2bar)
    [FiniteDimensional Q2 K]
    (localData : DyadicLocalData K)
    (branch : BranchArithmetic K)
    (cert : BranchCertificate localData branch) :
    Nonempty
      (ContinuousMulEquiv
        (candidateGroup localData.params branch.word)
        K.fixingSubgroup)
\end{verbatim}
The proof of this theorem should be the composition of a branch certificate
with a single generic reconstruction theorem.  It should not introduce an
isomorphism axiom for any field.

\subsection{Parameter and marking layer}

The first production modules should define:
\begin{verbatim}
inductive LabuteType
  | L
  | M (alpha : Nat)
  | N (alpha : Nat)

inductive Generator (n : Nat)
  | sigma
  | tau
  | wild (i : Fin (n + 1))

structure Marking (n : Nat) (G : Type*) where
  sigma : G
  tau : G
  wild : Fin (n + 1) -> G

structure FieldParameters where
  n : Nat
  residueDegree : Nat
  qK : Nat
  qK_eq : qK = 2 ^ residueDegree
\end{verbatim}
The semantic generator type is preferable to a bare
\(\operatorname{Fin}(n+3)\): it prevents coordinate-order mistakes in the
Fox complex while retaining the correct cardinality.  The old
four-generator marking should remain in place initially, with an adapter
identifying
\[
 (\sigma,\tau,x_0,x_1)
 \quad\longleftrightarrow\quad
 (\mathtt{sigma},\mathtt{tau},\mathtt{wild}\ 0,
  \mathtt{wild}\ 1)
\]
at \(n=1\).

The cyclotomic Frobenius data should be split into a canonical quotient
element and a chosen word representative:
\begin{verbatim}
structure CyclotomicFrobeniusDatum
    (C : Type*) [Group C] where
  r : Nat
  lambda : C ->* Multiplicative (ZMod (2 ^ r))
  lambda_surjective : Function.Surjective lambda

def inertiaImage (d : CyclotomicFrobeniusDatum C) : Subgroup C :=
  d.lambda.ker

def gammaCoset (d : CyclotomicFrobeniusDatum C) :
    C / d.inertiaImage := ...
\end{verbatim}
A representative of \(\mathtt{gammaCoset}\) should be chosen only when a
concrete word requires a Frobenius lift.  This avoids making the
presentation depend on a noncanonical representative.

\subsection{The four formal interfaces}

The mathematical B1--B4 split should be reflected literally in the Lean
API.

\paragraph{B1: \texttt{SourceBoundary} and \texttt{BoundaryPair}.}
Tame and maximal pro-\(2\) maps, equality of their full
\(\Z_2\)-characters, quotient identifications, joint surjectivity to the
fibre product, and the values of \(\nu\) on every marked generator.

\paragraph{B2: \texttt{ExactLiftingSemantics}.}
The maps
\[
 A\longrightarrow(\mathtt{Generator}\ n\to A)
 \longrightarrow A\times A,
\]
the chain identity, literal defect-change formula,
obstruction-vanishing equivalence, lift-torsor theorem, naturality, and
exactness in coefficients.

\paragraph{B3: \texttt{LinearDualityInterface}.}
The finite-word Stokes chain map, perfectness in degrees \(0/2\) and \(1\),
the \(H^2\)-invariant identification, adjoint connecting maps, and scalar
Hilbert normalization.  Cardinality formulas alone are insufficient.

\paragraph{B4: \texttt{QuadraticDeterminantInterface}.}
Candidate and local quadratic forms, nonsingularity, ramification kind,
the two Gauss formulas, and their equality.

The generic package should then read schematically:
\begin{verbatim}
structure SourceInterface
    (P : DyadicPresentation n)
    (K : IntermediateField Q2 Q2bar) where
  boundary : BoundaryPair n P.group K.fixingSubgroup
  exactSemantics : ExactLiftingSemantics P
  linear : LinearDualityInterface P
  quadratic : QuadraticDeterminantInterface P

theorem sourceInterface_contSurjCounts
    (I : SourceInterface P K) (G : Type*)
    [Group G] [Finite G] :
    Nat.card (ContSurj P.group G) =
      Nat.card (ContSurj K.fixingSubgroup G)

theorem sourceInterface_reconstruction
    (I : SourceInterface P K) :
    Nonempty (ContinuousMulEquiv P.group K.fixingSubgroup)
\end{verbatim}
The last theorem should call the existing profinite reconstruction API.
The boundary-framed induction should be generalized by replacing its
hard-coded source leaves with fields of
\texttt{Source}\-\texttt{Interface}; the
finite group recursion itself should remain unchanged.

\subsection{New foundational work}

The marked relative Labute theorem is the part least represented in the
current codebase.  It requires the following modules, in dependency order:
\begin{enumerate}
\item \texttt{FreeProTwo}: the maximal pro-\(2\) quotient of a free
profinite group, with transported generator and evaluation lemmas;
\item \texttt{Zassenhaus}: the filtration, commutator and square bounds,
and the graded pieces actually used by the proof;
\item \texttt{RestrictedLie}: the bracket, restricted square, and
degreewise derivations induced by Nielsen corrections;
\item \texttt{HallBasis}: only the basis elements and rank calculations
needed by the \(L\), \(M\), and \(N\) rows;
\item \texttt{LabuteWords}: the marked model relators, their abelian
vectors, initial forms, orientations, and unramified characters;
\item \texttt{RelativeLabute}: surjectivity of the compatible correction
map in the four graded rows \(L\), \(M_{\epsilon=0}\),
\(M_{\epsilon=1}\), and \(N\); and
\item \texttt{NielsenLimit}: convergence of corrections that are the
identity in successively deeper filtration quotients.
\end{enumerate}
The output must be an equivalence preserving both the full cyclotomic
orientation and the full \(\Z_2\)-valued unramified character.  A
mod-\(2\) compatibility theorem is not sufficient for B1.

\subsection{Local Gauss theorem}

The current local quadratic construction should first be parameterized
from the concrete source \(G_{\Q_2}\) to an arbitrary profinite group with
a \texttt{TateDualityG} instance.  The principal target is:
\begin{verbatim}
theorem localDeterminantGauss
    (localData : DyadicLocalData K)
    (D : TateDualityG K.fixingSubgroup 2)
    (dat : LocalDeterminantDatum K C V) :
    GeneralGaussDatum localData.params.n dat.d
      (H1 K.fixingSubgroup V)
\end{verbatim}
The output record contains quadraticity, nonsingularity,
\(\#H^1=2^{nd}\), evenness of \(nd\), ramification kind, and the formula
\[
 \Gsum(Q)=
 \begin{cases}
 (-1)^n2^{nd/2},&\text{unramified},\\
 2^{nd/2},&\text{ramified}.
 \end{cases}
\]
The unramified proof should expose the rank-\(n\) Hermitian decomposition
without specializing \(n\) early.  The ramified proof should be divided
into odd-layer vanishing, projectivity of the faithful image, construction
of the deep Lagrangian, and the normalized Evens--Kahn calculation.

\subsection{Implementation order and checkpoints}

The recommended merge order is:
\begin{enumerate}
\item parameters, semantic marking, and \(n=1\) adapters;
\item group-generic local cohomology and wrappers recovering the current
\(G_{\Q_2}\) declarations;
\item the general local Gauss theorem;
\item the \(n+3\)-generator Fox complex and exact lifting semantics;
\item relative Labute infrastructure and the marked pro-\(2\) theorem;
\item the source-interface records and the parameterized finite-target
induction;
\item the \(N_2\) instance as the first five-generator end-to-end test;
\item the remaining \(L\), \(N_\alpha\), and \(M_\alpha\) instances.
\end{enumerate}
At each checkpoint the existing \(G_{\Q_2}\) theorem should continue to
compile through compatibility wrappers.

\subsection{Axiom boundary and acceptance criteria}

The formalization should distinguish new mathematics from classical local
inputs.  Initially it is reasonable to bundle local reciprocity, the
oriented tame quotient, local Euler characteristic, Tate duality, and
finite-dyadic Hilbert-symbol facts.  The relative Labute project should
remove the current composite orientation axiom by deriving it from smaller
arithmetic facts; doing so need not immediately eliminate every
literature axiom.

Every merge checkpoint should enforce:
\begin{enumerate}
\item no \texttt{sorryAx} and no new undocumented axiom;
\item preservation of existing \(G_{\Q_2}\) theorem names through wrappers;
\item
\(\mathtt{generalGaussSign\ 1\ unramified=-1}\) and
\(\mathtt{generalGaussSign\ 2\ unramified=1}\);
\item equivalence of the \(n=1\) semantic marking and Fox complex with the
old \(\operatorname{Fin}4\) versions;
\item five generator coordinates and two relation coordinates in the
\(N_2\) complex;
\item preservation of the full \(\nu:\!-\to\Z_2\) and cyclotomic
orientation by every marked pro-\(2\) equivalence; and
\item use of \texttt{sourceInterface\_reconstruction} in the final field
theorems, with no field-specific presentation-isomorphism axiom.
\end{enumerate}

\subsection{Current scaffold}

The accompanying file
\path{work/execution/dyadic_general_skeleton.lean} now compiles against
the existing GQ2 Lake project.  It contains:
\begin{itemize}
\item the parameter, family, generator, and marking types;
\item functorial mapping of general markings;
\item the B1 boundary fibre product and two-source boundary pair;
\item the two-relator complex and explicit B2/B3 obligation records;
\item the arbitrary-dyadic Gauss-sign record and B4 comparison record;
\item the combined source-interface input record; and
\item an ordered milestone enumeration for project tooling.
\end{itemize}
It introduces no axiom and no \texttt{sorry}.  It remains a type-level
scaffold: the substantive relative Labute, local Gauss, Fox, and
source-interface proofs are the next implementation stages.

\section{Publication status and remaining work}\label{sec:status}

The present consolidation treats the following as standard literature
inputs:
\begin{enumerate}
\item local Euler--Poincar\'e characteristic and Tate duality
\cite[Chapter VII]{NSW};
\item the dyadic unit filtration, Hilbert orthogonality, and local
reciprocity \cite{SerreLocal};
\item the modular projectivity criterion for a module free over a Sylow
subgroup;
\item the normalized Evens--Kahn identity
\cite{Evens1963,Kahn1984,Kozlowski1984};
\item Labute's classification of Demu\v{s}kin groups
\cite{Labute1967}; and
\item Diekert's theorem in the unramified-\(i\) branch
\cite{Diekert1984}.
\end{enumerate}

Before publication, exact proposition numbers should replace the broad
citations for the square-class filtration and modular projectivity.
The component calculations should receive independent human review.
Finally, the existing Lean file is only an interface prototype; a
machine-checked proof would still need formalizations of the relative
Labute theorem, the local Gauss-sign theorem, the finite-word Stokes
construction, and the source-interface induction.

\begin{thebibliography}{99}

\bibitem{Brown1982}
K.~S. Brown,
\emph{Cohomology of Groups},
Graduate Texts in Mathematics 87, Springer, 1982.

\bibitem{Diekert1984}
V. Diekert,
\emph{\"Uber die absolute Galoisgruppe dyadischer Zahlk\"orper},
J. Reine Angew. Math. \textbf{350} (1984), 152--172.

\bibitem{DixonEtAl1982}
J.~D. Dixon, E.~W. Formanek, J.~C. Poland, and L. Ribes,
\emph{Profinite completions and isomorphic finite quotients},
J. Pure Appl. Algebra \textbf{23} (1982), 227--231.

\bibitem{Evens1963}
L. Evens,
\emph{A generalization of the transfer map in the cohomology of groups},
Trans. Amer. Math. Soc. \textbf{108} (1963), 54--65.

\bibitem{FriedJarden}
M.~D. Fried and M. Jarden,
\emph{Field Arithmetic}, fourth edition,
Ergebnisse der Mathematik und ihrer Grenzgebiete (3), vol.~11,
Springer, 2023.

\bibitem{JannsenWingberg1982}
U. Jannsen and K. Wingberg,
\emph{Die Struktur der absoluten Galoisgruppe \(p\)-adischer
Zahlk\"orper},
Invent. Math. \textbf{70} (1982), 71--98.

\bibitem{Kahn1984}
B. Kahn,
\emph{Classes de Stiefel--Whitney de formes quadratiques et de
repr\'esentations galoisiennes r\'eelles},
Invent. Math. \textbf{78} (1984), 223--256.

\bibitem{Koch2002}
H. Koch,
\emph{Galois Theory of \(p\)-Extensions},
Springer Monographs in Mathematics, Springer, 2002.

\bibitem{Kozlowski1984}
A. Kozlowski,
\emph{The Evens--Kahn formula for the total Stiefel--Whitney class},
Proc. Amer. Math. Soc. \textbf{91} (1984), 309--313.

\bibitem{Labute1967}
J.~P. Labute,
\emph{Classification of Demushkin groups},
Canad. J. Math. \textbf{19} (1967), 106--132.

\bibitem{MilneADT}
J.~S. Milne,
\emph{Arithmetic Duality Theorems}, second edition,
BookSurge, 2006.

\bibitem{NSW}
J. Neukirch, A. Schmidt, and K. Wingberg,
\emph{Cohomology of Number Fields}, second edition,
Springer, 2015.

\bibitem{RibesZalesskii}
L. Ribes and P. Zalesskii,
\emph{Profinite Groups}, second edition,
Springer, 2010.

\bibitem{RoeTurturean2026}
D. Roe and D. Turturean,
\emph{A presentation of the absolute Galois group of \(\Q_2\)},
preprint, 2026,
\url{https://roed314.github.io/gq2/paper/paper.html}.

\bibitem{Schneider2016}
F. Schneider,
\emph{On the structure of the absolute Galois group of local fields
with residue characteristic \(2\)},
Ph.D. thesis, Universit\"at Regensburg, 2016.

\bibitem{SerreLocal}
J.-P. Serre,
\emph{Local Fields},
Graduate Texts in Mathematics 67, Springer, 1979.

\bibitem{SerreProP}
J.-P. Serre,
\emph{Structure de certains pro-\(p\)-groupes},
S\'eminaire Bourbaki, vol.~8, expos\'e no.~252 (1962--1964), 145--155.

\bibitem{Zelvenskii1972}
I.~G. Zel'venski\u{\i},
\emph{The algebraic closure of a local field for \(p=2\)},
Izv. Akad. Nauk SSSR Ser. Mat. \textbf{36} (1972), 933--946;
English transl., Math. USSR-Izv. \textbf{6} (1972), 925--937.

\end{thebibliography}

\end{document}
